% Build from docs/writeups with xelatex milestone2.tex or pdflatex milestone2.tex.
\documentclass[11pt,letterpaper]{article}
\usepackage[letterpaper,left=.85in,right=.85in,top=.71in,bottom=.76in,footskip=.36in]{geometry}
\usepackage{iftex}
\ifPDFTeX\usepackage[T1]{fontenc}\usepackage{lmodern}\else\usepackage{fontspec}\fi
\usepackage{booktabs,graphicx,array,hyperref}
\hypersetup{hidelinks,pdftitle={Group 4 - Milestone 2 - Choice reversals in household lotteries},pdfauthor={Ishan Shivansh Bangroo, Zhi-Hong Jian, Tara Liu, Leona Pierce}}
\urlstyle{same}
\setlength{\parindent}{0pt}\setlength{\parskip}{4.5pt}
\newcommand{\head}[1]{\par\vspace{2pt}\noindent\textbf{#1}\hspace{.3em}}
\newcommand{\titleline}[1]{{\LARGE\bfseries #1\par}\vspace{3pt}}
\newcommand{\subline}[1]{{\large\bfseries #1\par}}
\newcommand{\analysispath}{output/analysis/milestone_2}
\IfFileExists{\analysispath/09_summary_6_blocks.tex}{}{\edef\analysispath{../../\analysispath}}
\begin{document}
{\small IST 597-005\quad Group 4\quad Milestone 2}\par
\titleline{Choice reversals in household lotteries}
Ishan Shivansh Bangroo\quad Zhi-Hong Jian\quad Tara Liu\quad Leona Pierce\footnote{\scriptsize\textbf{Contributions.} Leona collected the first M2 runs and suggested the counterfactual revision test. Zhi-Hong extended the runs, assembled the human benchmark and proposed the three analyses. Tara helped develop the baseline prompts and suggested total variation distance. Ishan specified the revision measures and prepared the comparison tables and report.}

\head{Question and sample.} Do models revise ticket allocations as people do once they learn that their tickets decide the donation? We compared 239 participants in Andreoni et al.\ (2020) with 880 model conversations across 11 settings. Each player divided ten tickets between two households. Before the draw, players could offset the computer's unequal split. After learning that their own tickets determined the winner, a 5/5 split gave each household an equal chance.

\head{Three comparisons.} Total variation distance (TVD) compares the five initial allocation categories. It is half the sum of absolute differences in category shares. Zero means an exact match. The six equally weighted blocks cover rounds 1--4 of the pure before-draw (4A) and after-reveal (4P) conditions, and rounds 5--8 of all four conditions. The 4A and 4P columns compare separate groups. Only observed choices enter the rates. The final three columns follow the same revision opportunities and weight each condition equally. ``Changed'' counts any change in allocation.

{\centering\small\setlength{\tabcolsep}{5.25pt}\renewcommand{\arraystretch}{1.03}
\def\plannedonly{}\input{\analysispath/09_summary_6_blocks.tex}\par}

\head{Findings.} Humans chose 5/5 on 23 of 284 initial choices in 4A and 62 of 144 in 4P. On eligible rounds, their 5/5 rate rose from 7.1\% to 56.4\% after the reveal. Eight of 11 settings chose 5/5 more often in 4P than in 4A. Eight increased after the reveal. Six showed both human patterns. Across six blocks, DeepSeek at high reasoning came closest on initial choices (mean TVD 0.315), yet its eligible 5/5 rate rose from 0\% to 87.5\%. All 102 Gemma changes with reasoning on ended at 5/5. Pooling the four human conditions gives 315 changes in 478 offers (65.9\%), with 243 ending at 5/5. These are descriptive comparisons of repeated choices.

\head{Interpretation.} Haiku with reasoning off kept all 160 eligible allocations. Only five started unequal. Holding the model and draws fixed, we could compare the original prompt, confirmation alone, and confirmation with the winning probabilities shown. Since no money changed hands, we could also test whether calling the donation hypothetical affects choices.

\head{Next step.} On a small subset, we will test new prompts, fuller descriptions of the lottery images, and temperature 2 separately. For the counterfactual test, we will show the model someone else's unequal split in a fresh chat. We will tell it that its group has the winning ticket and give it one chance to keep or change the split.

\newpage
{\small IST 597-005\quad Group 4\quad Milestone 2}\par
\titleline{Appendix A. Experiment and API setup}
\head{Design.} Rounds 1--4 were four A, four P, two A then two P, or two P then two A. In A the player allocated before learning which set held the winning ticket. In P a player made a choice only if their own set won. In rounds 5--8, all four allocations preceded the first reveal. Each player then had two opportunities to keep or change an allocation. Each conversation kept its own history. The computer gave household A 7, 2, 10, 1, 8, 3, 9 and 0 tickets in successive rounds. The code balanced wins within four-round blocks and seeded draws by run, condition and player number, sharing draws across model settings.

\head{Interface.} We called \url{https://openrouter.ai/api/v1/chat/completions}. The OpenRouter client reads \texttt{OPENROUTER\_API\_KEY} from the environment or local \texttt{.env} and sends it as a bearer credential. Requests supplied the model identifier, conversation messages and temperature 1.0. For reasoning, on/off supplied an enabled flag, medium/high supplied an effort level, and default omitted the override. These are requested settings. The saved histories do not establish which internal computation a provider used.

{\centering\small\setlength{\tabcolsep}{8pt}\renewcommand{\arraystretch}{1.08}
\begin{tabular}{@{}ll@{}}\toprule
Model identifier in the records & Reasoning requested\\\midrule
\texttt{anthropic/claude-haiku-5.5} & off, medium\\
\texttt{deepseek/deepseek-v4.1-flash-1017} & off, high\\
\texttt{google/gemini-3.8-flash} & medium\\
\texttt{google/gemma-4-31b-it} & off, on (earlier sample used default)\\
\texttt{openai/gpt-6-luna} & off, medium\\
\texttt{qwen/qwen3.8-flash} & off, medium\\\bottomrule
\end{tabular}\par}

\head{Records and checks.} Each main setting has 20 conversations per condition. Another 40 Gemma conversations used the default reasoning setting (10 per condition). Appendices B and C identify that earlier sample as Gemma default. The parser reads the last JSON object and accepts two whole, nonnegative counts summing to ten. The runner allows up to four attempts per question. Of the 8,246 accepted replies across all 920 records, 1,758 contained text outside the requested standalone object. All saved replies parsed to the recorded allocations. The full request-attempt log is unavailable.

\head{Human data.} We used ten de-identified sessions in the journal's replication archive. The four conditions contained 71, 72, 48 and 48 people respectively, in the order listed above. The survey recorded tickets for B, so we subtracted each observed count from ten to obtain tickets for A. Skipped P rounds had no participant allocation. We classified decisions as overcompensating, ex ante equalizing, mixed, 5/5 or reinforcing, relative to the computer's split.

\head{Reproduction.} From the repository root, install dependencies with \texttt{uv sync}, then run
\par{\small\ttfamily uv run python -m analysis.run\_analysis.milestone\_2\par}
The analysis reads \path{output/results/milestone_2/players.jsonl} and \path{human_benchmark/processed/players.jsonl}. It writes the figures, CSVs and LaTeX tables in \path{output/analysis/milestone_2/}. To collect a new sample, set a fresh \texttt{RUN\_NAME} in \path{main.py} and run \texttt{uv run python main.py} with a local credential. \path{prompt/prompt_a.py} holds the message templates.

{\small\head{Sources.} Andreoni, J., D. Aydin, B. Barton, B. D. Bernheim and J. Naecker (2020). When Fair Isn’t Fair: Understanding Choice Reversals Involving Social Preferences. \textit{Journal of Political Economy}, 128(5), 1673--1711. \href{https://doi.org/10.1086/705549}{doi:10.1086/705549}. \href{https://www.journals.uchicago.edu/doi/suppl/10.1086/705549/suppl_file/2016200data.zip}{Journal replication archive}. \href{https://github.com/leona5139/choice_reversals}{Group 4 project repository}.}

\newpage
{\small IST 597-005\quad Group 4\quad Milestone 2}\par
\titleline{Appendix B. Results by model setting}
The heatmap shows which conditions account for each model’s overall distance. In rounds 1--4, Haiku with reasoning off was closer to the human distribution in 4P (TVD 0.42) than in 4A (0.92). The connected markers show the direction of the change in 5/5 share, which the distance alone cannot show. In the human 4A condition, 97 of 142 eligible allocations changed (68.3\%), matching the published rate.

{\centering\includegraphics[width=\linewidth,height=.43\textheight,keepaspectratio]{\analysispath/06_tvd_6_blocks.png}\par}
\vspace{2pt}
{\centering\includegraphics[width=\linewidth,height=.39\textheight,keepaspectratio]{\analysispath/08_5_5_framing_and_revisions.png}\par}
\vspace{1pt}

\newpage
{\small IST 597-005\quad Group 4\quad Milestone 2}\par
\titleline{Appendix C. Twelve-block comparison}
Here the mean TVD gives equal weight to rounds 1--2, 3--4 and 5--8 in each of the four treatments. Qwen at medium reasoning has the lowest mean (0.336). The 5/5 and revision columns use the same decisions as the main table. The final row shows the earlier Gemma-default sample.

{\centering\small\setlength{\tabcolsep}{5.25pt}\renewcommand{\arraystretch}{1.03}
\input{\analysispath/09_summary_12_blocks.tex}\par}

\medskip
{\centering\includegraphics[width=\linewidth]{\analysispath/06_tvd_12_blocks.png}\par}

\end{document}
